\chapter{第八讲\quad 主观要件}
\label{lecture:08}
\noindent{\small\hyperref[ref:lecture:08]{查阅本讲扩展规则}}\par

主观要件必须绑定到同一个行为、同一时点、同一个待评结果或对象。一个人对甲的死亡可以有直接故意，对乙的死亡却只有过失。判定某项罪过以后，仍须与该罪的客观要件核对。

\section{第一节\quad 犯罪故意}

\paragraph{定义与法条}
刑法第十四条规定故意犯罪的认识与意志条件。令 $E_j$ 为罪名 $j$ 要求行为人认识的构成事实集合，$K_t(e)$ 表示行为时认识到事实 $e$ 或具备与该罪相应的可能性认识，$W_t(r)$ 表示希望待审结果，$L_t(r)$ 表示放任该结果。对本讲的典型故意结果犯，先按如下认识与意志结构核验
\[
\mathrm{Intent}_j(a,r)
  =\left(\bigwedge_{e\in E_j}K_{t_a}(e)\right)
     \land\bigl(W_{t_a}(r)\lor L_{t_a}(r)\bigr),
\]
其中 $t_a$ 为实行行为的时点。$E_j$ 可含对象种类、行为危险、特定身份等应认识的事实，按具体罪状取值。一般对象的姓名等非构成事实无需加入。明知事实的程度、规范性要素的外行理解和特定目的，要按分则核查，不能用公式代替条文判断。

\paragraph{推导}
甲把眼前站着的丙当作仇人乙，对这个人开枪并使其死亡。故意杀人罪要求认识对象是人，不以知道姓名为条件。题设给出 $K_{t_a}(\text{眼前对象是人})$、$W_{t_a}(\text{该人死亡})$、实际死亡和枪击归责。因此这一识别错误不使所需认识要素为假。已满足着手与其他归责条件的，依第二百三十二条的教材模型认定故意杀人既遂。若甲确实误将眼前的人当成猎物，$K_{t_a}(\text{对象是人})$ 不成立，不能从客观死亡倒推故意，另查过失。

\paragraph{时间一致性与边界}
令 $t_0<t_1<t_2$ 分别为谋划、实施和事后。$W_{t_0}(r)$ 或 $W_{t_2}(r)$ 单独均不能推出 $W_{t_1}(r)$。甲计划将来撞死乙，开车赴现场途中不慎撞死乙，须按实施碰撞时的罪过评价该碰撞。先前计划是否达到预备，另外判断。违反交通规则的故意同样不能推出对人身伤亡的故意。目的犯还须查目的与实行行为同时存在。先为私怨殴打，后才发现财物并起取财之意，不能以事后目的证明先前暴力已经是抢劫手段。

\paragraph{概括故意与择一故意}
用 $\Omega$ 表示行为人认识到的结果情形。概括故意常见于对多名对象投放危险，行为人希望或接受可能发生若干死亡，却未特定名单和数量。择一故意常见于行为人按其计划认为结果将在若干目标中择一发生，并对各备选结果都抱有相应故意。二者讨论故意涉及哪些对象或结果。直接、间接故意则讨论意志态度。逐个结果判断实际发生与故意范围，不能因只实施一次行为便只审查一个对象，也不能在欠缺实际危险或着手时直接认定对未死亡对象成立未遂。

\section{第二节\quad 罪过形式的区分}

\paragraph{定义与规则层级}
令 $F$ 表示行为时已预见该结果有现实发生可能，$J$ 表示应当且能够预见，$A$ 表示依当时可行且应当采用的措施可以避免结果。$W$ 为希望，$L$ 为放任，$T$ 为轻信可避免。第十四条至第十六条是法条。下面是按教材拆分的判断核心，每行都假定已按具体罪名审查所需客观风险与结果。
\[
\begin{array}{c|l}
\text{罪过或排除情形} & \text{核心辨别条件}\\ \hline
\text{直接故意} & \text{明知必然或可能发生，并希望发生}\\
\text{间接故意} & F\land L\\
\text{过于自信的过失} & F\land T\land A\\
\text{疏忽大意的过失} & \neg F\land J\land A\\
\text{意外事件} & \neg F\land\neg J\text{，且无故意或过失}\\
\text{不可抗力} & \neg A\text{，且损害因不能抗拒的原因发生、无罪过}
\end{array}
\]
第十五条还要求有法定过失处罚条款。教材将明知伴随结果必然发生而仍实行的典型案例按直接故意评价。这里另行评价行为的意志，不能简单按口头「不希望」认作过失。意外事件与不可抗力可按不可预见和不可抗拒区分。不可抗力在法条上不以已经预见为必要条件。

\paragraph{推导，同一行为对两个结果分别求值}
甲为杀乙投毒，已想到共同生活的丙可能误食。甲确实不愿丙受害，采取自以为足够的隔离措施，但依题设该措施明显不足，另有可行措施可以避免丙误食。乙、丙均因该投毒死亡。对乙，$K_{t_a}$ 与 $W_{t_a}$ 均满足，对其死亡是直接故意。对丙，题设支持
\[
F_{\text{丙}}=1,\quad T_{\text{丙}}=1,\quad
A_{\text{丙}}=1,\quad L_{\text{丙}}=0
   \ \Longrightarrow\ \mathrm{ConsciousNeg}_{\text{丙}}.
\]
按已证死亡归责及第十五条、第二百三十三条审查过失致人死亡，最后再审一行为触及多个罪状的竞合规则。若题设改为甲已预见丙可能误食，并以「丙吃了也无所谓」接受该风险，则变化的是 $T,L$，丙的死亡改按间接故意审查。

\paragraph{边界}
是否采取措施只是证据线索。某人没有采取措施，可能轻信自己的经验，不能单由此推出放任。形式上摆出一项措施，也可能明知措施几乎不起作用而接受死亡。须结合危险程度、能力、具体措施、行为前后的言行和其他证据。若早先能够且应当避免，不能只截取事故发生的最后一秒说无法抗拒。若所有可行方案当时都无法阻止，且无先前制造风险的过错，则不能以「已经预见」直接认定有过失。若 $\Gamma$ 既不能证明 $F$，也不能证明 $\neg F$，应保留待证，不能自动落入疏忽大意。

\section{第三节\quad 事实认识错误}

\paragraph{统一检查式}
先定义三个角色。$p$ 是计划所指对象，$a$ 是实际瞄准或接触的对象，$v$ 是实际被害对象。$\tau(z)$ 表示法律上对对象、结果和行为方式的构成要件分类。设 $G(b,r)$ 是第七讲已经证成的实行行为 $b$ 与结果 $r$ 的归责关系。按次序先查是否着手、$G$ 是否成立，再查主客观的 $\tau$ 是否一致，最后区分对象、打击、因果流程，决定是否需要展示不同解释规则。

\paragraph{同一构成要件内，对象与打击的推导}
当 $a=v$ 而行为人把 $a$ 的身份误认成 $p$，通常是对象错误。甲瞄准的是眼前的丙，并打中丙。即使把丙叫作乙，只要所需的对象类别是人，丙仍是故意指向的实际对象。若实际瞄准乙且偏离击中丙，则 $a\ne v$，进入打击错误。题设假定甲欲杀乙已着手，丙因偏弹死亡，甲没有接受丙死亡但具备对丙的过失。记
\[
\Gamma\vdash I_{\text{乙}}\land\neg G(b,r_{\text{乙}})
               \land G(b,r_{\text{丙}})\land N_{\text{丙}},
\]
其中 $I$ 表示杀人故意，$N$ 表示已证明的过失。两种教材争议规则分别推导：
\begin{enumerate}
\item 具体符合说强调实际指向。对乙按故意杀人未遂判断。对丙以 $G(b,r_{\text{丙}})\land N_{\text{丙}}$ 接入过失致人死亡，再处理想象竞合。若不能证明对丙的过失，不成立这一过失罪。
\item 法定符合说在本案同一构成要件的范围内评价偏离，对丙按故意杀人既遂，对乙按未遂，并处理同一行为的竞合。这里使用的是学说规则，不能从 $\exists y\,I(y)$ 以谓词逻辑推出 $I(\text{丙})$，也不等于证明甲在心理事实上希望丙死亡。
\end{enumerate}
若甲其实在射击时也放任旁边丙死亡，则依直接已证的主观事实评判丙的死亡，不需要靠打击错误规则补足。特殊对象的法定身份本身属于构成要件时，身份误认是否重要，应核对该要件有无实现。对象错误和打击错误可以同题出现，应画出 $p,a,v$ 的对应关系逐项分析。

\paragraph{同一构成要件内，因果流程错误}
若同一行为已经着手并具有杀人故意，推乙下井，本想淹死乙，乙因该次坠落撞击死亡。题设给出 $G(b,r)$ 且死亡实现了推落产生的危险。若流程差异不重要，教材认定故意杀人既遂。若同一时间另有独立因素单独造成死亡，$G$ 不成立，即便死的是预想对象，也不能按原行为既遂。

结果推迟时，写明两个行为 $b_1,b_2$。甲先故意击昏乙，误以为乙死，后来抛入水中使其溺亡。令 $G_k=G(b_k,r)$。给定 $G_2=1$ 且第二行为有死亡过失，教材展示的争议为：
\begin{enumerate}
\item 若认为抛放在具体情境中由先前伤害所引发，可推 $G_1=1$。
\item 若认为把活人误作死者是独立异常介入，继续比较两段机制。若先前失去意识和后续入水均实质参与，则 $G_1=1$。若只有后行为风险决定死亡，则 $G_1=0$。
\end{enumerate}
在 $G_1=1$ 的分支，第一行为按故意杀人既遂评价，教材以避免重复评价同一死亡结果处理与第二行为过失致死的关系。在 $G_1=0$ 的分支，题设满足条件时是故意杀人未遂和另行的过失致人死亡，按两个独立行为处理。若被害人只是佯装死亡，前行为未造成昏迷，不能直接沿用第一分支所凭事实。

结果提前时，先审先行行为是否已经对法益造成现实、紧迫、直接的危险，记作 $S_1$。若 $S_1=0$，仅在计划的预备阶段发生死亡，不能以结果补足故意罪既遂的着手。根据具体情节可能同时评价犯罪预备和过失致死。若 $S_1=1$，例如按杀人计划先扼颈控制、再打算抛水，扼颈已经使人死亡，教材的多数说在归责及主观要件成立时按既遂。少数说认为该死亡仅可按过失归责，分别审查故意未遂和过失致死。可见时间早晚本身并不决定阶段，关键在前行为是否已经是实行行为及其故意范围。

\paragraph{不同构成要件间，抽象的认识错误}
以 $O_j$ 表示罪名 $j$ 的已证客观条件，$I_j$ 表示已证主观条件。对不同构成要件，逐罪求交。甲故意射击贵重器物，已经着手，但偏弹击毙人，题设确认没有杀人故意且有死亡过失。器物未毁意味着缺少故意毁坏财物既遂的结果。人的死亡属于另一构成，按过失致人死亡的条件审查。前罪是否处罚未遂，还要达到该罪要求的法定程度。不能把死亡认作毁物，也不能把毁物故意直接当作杀人故意。

若教材主张两个罪状在个案中有包容关系，需证明两类要件存在可投影的共同部分。用 $P$ 表示普通财物的要件，$Q$ 表示枪支特征。甲为盗取普通财物取走包，包内实为枪支。客观上满足 $P\land Q$，主观上只认识 $P$。从 $P\land Q$ 可推出 $P$，从 $I(P)$ 不能推出 $I(Q)$。若盗窃罪的数额或法定行为方式等其余条件满足，可按普通盗窃既遂。未证明认识枪支特征，不能成立盗窃枪支的故意。包容关系须对照具体条文与证据，不可按罪名轻重任意建立。

\section*{闭卷记忆与自测}

\paragraph{复述路线}
拿笔写「谁、何时、知什么、想怎样」。看到罪过，先按结果分行，然后用「已见看态度，未见看应见，过失问能避」填 $F,W,L,T,J,A$。看到错误，按「着手、归责、同异、偏差、竞合」写五步。辨对象和打击时在纸上写 $p,a,v$，分别记录计划、实际瞄准和实际受害者。背完主流结论后，遮住答案，只改变一项事实重新推出结果。隔一天再用同一张五步表做一题。

\paragraph{题一}
甲预见操作会引发伤亡，以为普通防护足够，实际上有合规措施可避免。若改成甲认为伤亡发生也无所谓，哪些项变化？

\paragraph{核对}
两案均有 $F$。首案有 $T,A$，按过于自信的过失核查。变式用 $L$ 取代 $T$，按间接故意核查。结论还要以所审罪状的客观事实成立为前提。

\paragraph{题二}
「把丙当乙，瞄准丙并击中丙」和「瞄准乙，偏弹击中丙」最先不同的是哪一项？

\paragraph{核对}
首案 $a=v=\text{丙}$，是对象错误。次案 $a=\text{乙}\ne v$，是打击错误。两案对象都属于人时，仍需分别处理，不能只按「死者非原计划人」合并为一类。

\paragraph{题三}
实行前有杀人计划，实际在预备阶段造成死亡，能否仅因实际死亡认定故意杀人既遂？

\paragraph{核对}
先审着手 $S_1$。题设为真预备则 $S_1=0$，不能以死亡补上着手。再分开论预备、死亡过失与竞合。有明确法定特殊规则时，应另行适用。
